\section{Introduction}

Although stock-market turnover has long been a reference point for practitioner financial analysts, the academic asset-pricing literature has predominantly focused on explaining asset prices, leaving trading volume data unexplained. The strong emphasis on prices rather than quantities stems from a fundamental principle in finance: when arbitrage opportunities are minimal, all asset prices can be explained by a common discount factor. Unsurprisingly, a guiding principle of the theoretical asset-pricing literature has been that by understanding the drivers of this discount factor, we can identify relevant patterns for real-world investors, such as excess returns, predictability, and relative volatility, without developing predictions about volume.

Grounded on investor-level data, behavioral finance attempts to link asset prices, and consequently the stochastic discount factor, to investor irrationality and heterogeneity. Notable contributions include the work of Robert Shiller and Andrei Shleifer\todo{some cite}, who have investigated how investor sentiment and cognitive biases can impact asset prices. Yet, Harris and Raviv (1993) noted that divergent beliefs drive asset prices and trading volume. Thus, studying how belief heterogeneity explains trading volumes is helpful to derive a consistent set of predictions that bridge the micro evidence of investor beliefs with the macro phenomena observed in asset prices.

However, due to the emphasis on prices, we lack transparent formulas linking volume to different belief heterogeneity. This paper turns in that due homework. In this study, we derive explicit formulas for stock-market turnover as a function of investor belief heterogeneity. We do so in the simplest model possible: we study volume in a complete-markets model of log-utility investors who hold heterogeneous beliefs about fundamentals. We analyze how various belief models influence the correlation between stock-market turnover, asset prices, and business cycles. Among these models, we find that only extrapolative beliefs align closely with empirical data, suggesting that investors' tendency to project past trends into the future significantly drives market turnover.

A third property regards stock-market turnover, a feature that receives little attention in the macro-finance literature. We show that when some households are extrapolative, there is a positive correlation between stock-market turnover and measured disagreement. Only for this class of beliefs is turnover particularly large when the economy enters a recession, a property that we verify holds in the data. 



